\documentclass[11pt,a4paper]{article}
\usepackage[margin=0.7in]{geometry}
\usepackage[T1]{fontenc}
\usepackage[utf8]{inputenc}
\usepackage{inconsolata}
\usepackage{xcolor}
\usepackage{listings}
\usepackage{hyperref}
\usepackage{titlesec}
\usepackage{longtable}

\definecolor{codebg}{RGB}{248,248,248}
\definecolor{codeframe}{RGB}{220,220,220}
\definecolor{codetitle}{RGB}{25,63,108}

\lstdefinestyle{labcode}{
  basicstyle=\ttfamily\small,
  backgroundcolor=\color{codebg},
  frame=single,
  rulecolor=\color{codeframe},
  frameround=tttt,
  breaklines=true,
  breakatwhitespace=false,
  columns=fullflexible,
  keepspaces=true,
  showstringspaces=false,
  tabsize=4,
  numbers=left,
  numberstyle=\tiny\color{gray},
  numbersep=8pt,
  xleftmargin=8pt,
  xrightmargin=8pt
}

\titleformat{\section}{\Large\bfseries\color{codetitle}}{\thesection}{0.5em}{}
\titleformat{\subsection}{\large\bfseries}{\thesubsection}{0.5em}{}

\title{Network Lab Final Cheat Sheet\\\large Complete Source Dump (All Attached Folders)}
\author{}
\date{\today}

\begin{document}
\maketitle
\tableofcontents
\vspace{0.5em}

\section{Folder: Lab6}

\subsection*{Lab06Network.ned}
\begin{lstlisting}[style=labcode]
package lab06_network;

import inet.node.ethernet.Eth100M;
import inet.node.ethernet.EthernetSwitch;
import inet.node.inet.StandardHost;
import inet.networklayer.configurator.ipv4.Ipv4NetworkConfigurator;

network Lab06Network
{
    @display("bgb=916,555");
    submodules:
        configurator: Ipv4NetworkConfigurator {
            @display("p=50,50");
        }
        // Lab 1
        lab1_hostA: StandardHost {
            @display("p=100,100");
        }
        lab1_hostB: StandardHost {
            @display("p=100,200");
        }
        lab1_switch: EthernetSwitch {
            @display("p=200,150");
        }

        // Lab 2
        lab2_hostA: StandardHost {
            @display("p=300,100");
        }
        lab2_hostB: StandardHost {
            @display("p=381,54");
        }
        lab2_switch: EthernetSwitch {
            @display("p=336,161");
        }

        // Lab 3
        lab3_hostA: StandardHost {
            @display("p=500,100");
        }
        lab3_hostB: StandardHost {
            @display("p=584,53");
        }
        lab3_switch: EthernetSwitch {
            @display("p=489,161");
        }

        // Lab 4
        lab4_hostA: StandardHost {
            @display("p=700,100");
        }
        lab4_hostB: StandardHost {
            @display("p=866,99");
        }
        lab4_switch: EthernetSwitch {
            @display("p=737,216");
        }

        // Lab 5
        lab5_hostA: StandardHost {
            @display("p=599,358");
        }
        lab5_hostB: StandardHost {
            @display("p=610,437");
        }
        lab5_switch: EthernetSwitch {
            @display("p=500,400");
        }

        // Core switch to connect all labs
        coreSwitch: EthernetSwitch {
            @display("p=164,357");
        }

    connections:
        // Lab 1
        lab1_hostA.ethg++ <--> Eth100M <--> lab1_switch.ethg++;
        lab1_hostB.ethg++ <--> Eth100M <--> lab1_switch.ethg++;
        lab1_switch.ethg++ <--> Eth100M <--> coreSwitch.ethg++;

        // Lab 2
        lab2_hostA.ethg++ <--> Eth100M <--> lab2_switch.ethg++;
        lab2_hostB.ethg++ <--> Eth100M <--> lab2_switch.ethg++;
        lab2_switch.ethg++ <--> Eth100M <--> coreSwitch.ethg++;

        // Lab 3
        lab3_hostA.ethg++ <--> Eth100M <--> lab3_switch.ethg++;
        lab3_hostB.ethg++ <--> Eth100M <--> lab3_switch.ethg++;
        lab3_switch.ethg++ <--> Eth100M <--> coreSwitch.ethg++;

        // Lab 4
        lab4_hostA.ethg++ <--> Eth100M <--> lab4_switch.ethg++;
        lab4_hostB.ethg++ <--> Eth100M <--> lab4_switch.ethg++;
        lab4_switch.ethg++ <--> Eth100M <--> coreSwitch.ethg++;

        // Lab 5
        lab5_hostA.ethg++ <--> Eth100M <--> lab5_switch.ethg++;
        lab5_hostB.ethg++ <--> Eth100M <--> lab5_switch.ethg++;
        lab5_switch.ethg++ <--> Eth100M <--> coreSwitch.ethg++;
}
\end{lstlisting}

\subsection*{omnetpp.ini}
\begin{lstlisting}[style=labcode]
[General]
network = Lab06Network
sim-time-limit = 15s


*.lab*_host*.numApps = 1
*.lab*_host*.app[0].typename = "PingApp"
*.lab*_host*.app[0].packetSize = 500B
*.lab*_host*.app[0].sendInterval = 1s

*.*.ipv4.arp.typename = "GlobalArp"


*.lab1_hostA.app[0].destAddr = "lab1_hostB"
*.lab2_hostA.app[0].destAddr = "lab2_hostB"
*.lab3_hostA.app[0].destAddr = "lab3_hostB"
*.lab4_hostA.app[0].destAddr = "lab4_hostB"
*.lab5_hostA.app[0].destAddr = "lab5_hostB"
\end{lstlisting}

\section{Folder: Lab6Practice}

\subsection*{lab6.ned}
\begin{lstlisting}[style=labcode]
package lab6practice;

import inet.networklayer.configurator.ipv4.Ipv4NetworkConfigurator;
import inet.node.ethernet.EthernetSwitch;
import inet.node.inet.StandardHost;
import inet.node.ethernet.Eth100M;
import ned.IdealChannel;


network lab6Network
{
    @display("bgb=1171,573");
    submodules:
        configurator: Ipv4NetworkConfigurator {
            @display("p=96,76");
        }
        lab1HostA: StandardHost {
            @display("p=637,62");
        }
        lab1HostB: StandardHost {
            @display("p=803,84");
        }
        lab2HostA: StandardHost {
            @display("p=1042,185");
        }
        lab2HostB: StandardHost {
            @display("p=1049,282");
        }
        lab3HostA: StandardHost {
            @display("p=628,507");
        }
        lab3HostB: StandardHost {
            @display("p=488,515");
        }
        lab4HostB: StandardHost {
            @display("p=130,298");
        }
        lab4HostA: StandardHost {
            @display("p=145,394");
        }
        lab5HostB: StandardHost {
            @display("p=278,119");
        }
        lab5HostA: StandardHost {
            @display("p=457,62");
        }
        lab5Switch: EthernetSwitch {
            @display("p=400,165");
        }
        lab1Switch: EthernetSwitch {
            @display("p=659,152");
        }
        lab2Switch: EthernetSwitch {
            @display("p=862,281");
        }
        lab3Switch: EthernetSwitch {
            @display("p=527,433");
        }
        lab4Switch: EthernetSwitch {
            @display("p=266,350");
        }
        coreSwitch: EthernetSwitch {
            @display("p=581,291");
        }
    connections:
        lab5HostB.ethg++ <--> Eth100M <--> lab5Switch.ethg++;
        lab5HostA.ethg++ <--> Eth100M <--> lab5Switch.ethg++;
        lab1HostA.ethg++ <--> Eth100M <--> lab1Switch.ethg++;
        lab1HostB.ethg++ <--> Eth100M <--> lab1Switch.ethg++;
        lab2HostA.ethg++ <--> Eth100M <--> lab2Switch.ethg++;
        lab2Switch.ethg++ <--> Eth100M <--> lab2HostB.ethg++;
        lab3Switch.ethg++ <--> Eth100M <--> lab3HostB.ethg++;
        lab3Switch.ethg++ <--> Eth100M <--> lab3HostA.ethg++;
        lab4Switch.ethg++ <--> Eth100M <--> lab4HostA.ethg++;
        lab4HostB.ethg++ <--> Eth100M <--> lab4Switch.ethg++;
        lab4Switch.ethg++ <--> Eth100M <--> coreSwitch.ethg++;
        lab5Switch.ethg++ <--> Eth100M <--> coreSwitch.ethg++;
        lab1Switch.ethg++ <--> Eth100M <--> coreSwitch.ethg++;
        lab2Switch.ethg++ <--> Eth100M <--> coreSwitch.ethg++;
        lab3Switch.ethg++ <--> Eth100M <--> coreSwitch.ethg++;
}
\end{lstlisting}

\subsection*{omnetpp.ini}
\begin{lstlisting}[style=labcode]
[General]
network = lab6Network
sim-time-limit = 15s


*.lab1HostA.numApps = 3

*.lab1HostA.app[0].typename = "PingApp"
*.lab1HostA.app[0].destAddr = "lab2HostB"
*.lab1HostA.app[0].packetSize = 500B
*.lab1HostA.app[0].sendInterval = 1s

*.lab1HostA.app[1].typename = "PingApp"
*.lab1HostA.app[1].destAddr = "lab3HostB"
*.lab1HostA.app[1].packetSize = 500B
*.lab1HostA.app[1].sendInterval = 1s

*.lab1HostA.app[2].typename = "PingApp"
*.lab1HostA.app[2].destAddr = "lab4HostB"
*.lab1HostA.app[2].packetSize = 500B
*.lab1HostA.app[2].sendInterval = 1s
\end{lstlisting}

\subsection*{package.ned}
\begin{lstlisting}[style=labcode]
package lab6practice;

@license(LGPL);
\end{lstlisting}

\section{Folder: Lab7Network}

\subsection*{ip.xml}
\begin{lstlisting}[style=labcode]
<?xml version="1.0" encoding="UTF-8"?>
<config>
	<!-- engineer -->
	<interface hosts="centralRouter" names="eth0" address="192.35.8.1" netmask="255.255.255.192"/>
<!-- 	Market -->
	<interface hosts="centralRouter" names="eth1" address="192.35.8.129" netmask="255.255.255.240"/>
<!-- 	Sales -->
	<interface hosts="centralRouter" names="eth2" address="192.35.8.97" netmask="255.255.255.240"/>
<!--  	Admin -->
 	<interface hosts="centralRouter" names="eth3" address="192.35.8.113" netmask="255.255.255.240"/>
<!-- 	Management -->
	<interface hosts="centralRouter" names="eth4" address="192.35.8.65" netmask="255.255.255.224"/>

<interface hosts="engHost1" address="192.35.8.10" netmask="255.255.255.192"/>
    <interface hosts="engHost2" address="192.35.8.11" netmask="255.255.255.192"/>
    <interface hosts="engHost3" address="192.35.8.12" netmask="255.255.255.192"/>
    
    
   <interface hosts="mktHost1" address="192.35.8.130" netmask="255.255.255.240"/>
    <interface hosts="mktHost2" address="192.35.8.131" netmask="255.255.255.240"/>
    <interface hosts="mktHost3" address="192.35.8.132" netmask="255.255.255.240"/>
    
       
 <interface hosts="salesHost1" address="192.35.8.98" netmask="255.255.255.240"/>
    <interface hosts="salesHost2" address="192.35.8.99" netmask="255.255.255.240"/>
    <interface hosts="salesHost3" address="192.35.8.100" netmask="255.255.255.240"/>


  <interface hosts="adminHost1" address="192.35.8.114" netmask="255.255.255.240"/>
    <interface hosts="adminHost2" address="192.35.8.115" netmask="255.255.255.240"/>
    <interface hosts="adminHost3" address="192.35.8.116" netmask="255.255.255.240"/>
    
      <interface hosts="mgmtHost1" address="192.35.8.66" netmask="255.255.255.224"/>
    <interface hosts="mgmtHost2" address="192.35.8.67" netmask="255.255.255.224"/>
    <interface hosts="mgmtHost3" address="192.35.8.68" netmask="255.255.255.224"/>


	<route hosts="centralRouter" destination="192.35.8.0" netmask="255.255.255.192" interface="eth0"/>
    <route hosts="centralRouter" destination="192.35.8.128" netmask="255.255.255.240"   interface="eth1"/>
    <route hosts="centralRouter" destination="192.35.8.96" netmask="255.255.255.240" interface="eth2"/>
    <route hosts="centralRouter" destination="192.35.8.112" netmask="255.255.255.240" interface="eth3"/>
    <route hosts="centralRouter" destination="192.35.8.64" netmask="255.255.255.224"  interface="eth4"/>

	<route hosts="engHost*" destination="*" netmask="*" gateway="192.35.8.1"/>
    <route hosts="mktHost*" destination="*" netmask="*" gateway="192.35.8.129"/>
    <route hosts="salesHost*" destination="*" netmask="*" gateway="192.35.8.97"/>
    <route hosts="adminHost*" destination="*" netmask="*" gateway="192.35.8.113"/>
    <route hosts="mgmtHost*" destination="*" netmask="*" gateway="192.35.8.65"/>


</config>
\end{lstlisting}

\subsection*{Lab07Network.ned}
\begin{lstlisting}[style=labcode]
package lab07network;

import inet.node.ethernet.Eth100M;
import inet.node.ethernet.EthernetSwitch;
import inet.node.inet.StandardHost;
import inet.node.inet.Router;
import inet.networklayer.configurator.ipv4.Ipv4NetworkConfigurator;

network Lab07Network
{
    @display("bgb=1200,800");

    submodules:
        configurator: Ipv4NetworkConfigurator {
            @display("p=50,50");
            addStaticRoutes = false;
        }

        centralRouter: Router {
            @display("p=600,343;i=abstract/router");
        }

        engSwitch: EthernetSwitch {
            @display("p=300,200");
        }
        engHost1: StandardHost {
            @display("p=127,190");
        }
        engHost2: StandardHost {
            @display("p=236,100");
        }
        engHost3: StandardHost {
            @display("p=489,137");
        }

        mktSwitch: EthernetSwitch {
            @display("p=900,200");
        }
        mktHost1: StandardHost {
            @display("p=743,50");
        }
        mktHost2: StandardHost {
            @display("p=941,50");
        }
        mktHost3: StandardHost {
            @display("p=1098,127");
        }


        salesSwitch: EthernetSwitch {
            @display("p=200,500");
        }
        salesHost1: StandardHost {
            @display("p=80,535");
        }
        salesHost2: StandardHost {
            @display("p=248,622");
        }
        salesHost3: StandardHost {
            @display("p=362,537");
        }

        adminSwitch: EthernetSwitch {
            @display("p=600,600");
        }
        adminHost1: StandardHost {
            @display("p=442,711");
        }
        adminHost2: StandardHost {
            @display("p=743,722");
        }
        adminHost3: StandardHost {
            @display("p=756,590");
        }


        mgmtSwitch: EthernetSwitch {
            @display("p=899,397");
        }
        mgmtHost1: StandardHost {
            @display("p=986,561");
        }
        mgmtHost2: StandardHost {
            @display("p=1121,500");
        }
        mgmtHost3: StandardHost {
            @display("p=1089,343");
        }

    connections:

        engHost1.ethg++ <--> Eth100M <--> engSwitch.ethg++;
        engHost2.ethg++ <--> Eth100M <--> engSwitch.ethg++;
        engHost3.ethg++ <--> Eth100M <--> engSwitch.ethg++;
        engSwitch.ethg++ <--> Eth100M <--> centralRouter.ethg++;
        mktHost1.ethg++ <--> Eth100M <--> mktSwitch.ethg++;
        mktHost2.ethg++ <--> Eth100M <--> mktSwitch.ethg++;
        mktHost3.ethg++ <--> Eth100M <--> mktSwitch.ethg++;
        mktSwitch.ethg++ <--> Eth100M <--> centralRouter.ethg++;
        salesHost1.ethg++ <--> Eth100M <--> salesSwitch.ethg++;
        salesHost2.ethg++ <--> Eth100M <--> salesSwitch.ethg++;
        salesHost3.ethg++ <--> Eth100M <--> salesSwitch.ethg++;
        salesSwitch.ethg++ <--> Eth100M <--> centralRouter.ethg++;
        adminHost1.ethg++ <--> Eth100M <--> adminSwitch.ethg++;
        adminHost2.ethg++ <--> Eth100M <--> adminSwitch.ethg++;
        adminHost3.ethg++ <--> Eth100M <--> adminSwitch.ethg++;
        adminSwitch.ethg++ <--> Eth100M <--> centralRouter.ethg++;
        mgmtHost1.ethg++ <--> Eth100M <--> mgmtSwitch.ethg++;
        mgmtHost2.ethg++ <--> Eth100M <--> mgmtSwitch.ethg++;
        mgmtHost3.ethg++ <--> Eth100M <--> mgmtSwitch.ethg++;
        mgmtSwitch.ethg++ <--> Eth100M <--> centralRouter.ethg++;
}
\end{lstlisting}

\subsection*{omnetpp.ini}
\begin{lstlisting}[style=labcode]
[General]
network = Lab07Network
sim-time-limit = 30s

*.configurator.config = xmldoc("ip.xml")
*.configurator.assignAdress = true
*.configurator.assignRoutes = false
*.configurator.addStaticRoutes = false
*.configurator.addDefaultRoutes = false

*.engHost1.numApps = 1
*.engHost1.app[0].typename = "PingApp"
*.engHost1.app[0].destAddr = "mgmtHost2"
\end{lstlisting}

\section{Folder: lab7Practice}

\subsection*{ip2.xml}
\begin{lstlisting}[style=labcode]
<?xml version="1.0" encoding="UTF-8"?>
<config>

    <!-- ===== ROUTER A INTERFACES ===== -->
    <interface hosts="routerA" names="eth3" address="192.35.8.161" netmask="255.255.255.252"/>
    <interface hosts="routerA" names="eth0" address="192.35.8.1"   netmask="255.255.255.192"/>
    <interface hosts="routerA" names="eth1" address="192.35.8.145" netmask="255.255.255.240"/>
    <interface hosts="routerA" names="eth2" address="192.35.8.97"  netmask="255.255.255.224"/>

    <!-- ===== ROUTER B INTERFACES ===== -->
    <interface hosts="routerB" names="eth2" address="192.35.8.162" netmask="255.255.255.252"/>
    <interface hosts="routerB" names="eth0" address="192.35.8.65"  netmask="255.255.255.224"/>
    <interface hosts="routerB" names="eth1" address="192.35.8.129" netmask="255.255.255.240"/>

    <!-- ===== HOST INTERFACES ===== -->
    <interface hosts="engHostA"  address="192.35.8.2" netmask="255.255.255.192"/>
    <interface hosts="engHostB"  address="192.35.8.3" netmask="255.255.255.192"/>
    <interface hosts="engHostC"  address="192.35.8.4" netmask="255.255.255.192"/>

    <interface hosts="marketingHostA" address="192.35.8.146" netmask="255.255.255.240"/>
    <interface hosts="marketingHostB" address="192.35.8.147" netmask="255.255.255.240"/>
    <interface hosts="marketingHostC" address="192.35.8.148" netmask="255.255.255.240"/>

    <interface hosts="salesHostA" address="192.35.8.98"  netmask="255.255.255.224"/>
    <interface hosts="salesHostB" address="192.35.8.99"  netmask="255.255.255.224"/>
    <interface hosts="salesHostC" address="192.35.8.100" netmask="255.255.255.224"/>

    <interface hosts="mgmtHostA" address="192.35.8.66" netmask="255.255.255.224"/>
    <interface hosts="mgmtHostB" address="192.35.8.67" netmask="255.255.255.224"/>
    <interface hosts="mgmtHostC" address="192.35.8.68" netmask="255.255.255.224"/>

    <interface hosts="adminHostA" address="192.35.8.130" netmask="255.255.255.240"/>
    <interface hosts="adminHostB" address="192.35.8.131" netmask="255.255.255.240"/>
    <interface hosts="adminHostC" address="192.35.8.132" netmask="255.255.255.240"/>

    <!-- ===== STATIC ROUTES  ROUTERS ===== -->
    <!-- Router A: reach Router B's networks (Mgmt + Admin) via RouterB's link IP -->
    <route hosts="routerA" destination="192.35.8.64"  netmask="255.255.255.224" gateway="192.35.8.162"/>
    <route hosts="routerA" destination="192.35.8.128" netmask="255.255.255.240" gateway="192.35.8.162"/>

    <!-- Router B: reach Router A's networks (Eng + Marketing + Sales) via RouterA's link IP -->
    <route hosts="routerB" destination="192.35.8.0"   netmask="255.255.255.192" gateway="192.35.8.161"/>
    <route hosts="routerB" destination="192.35.8.144" netmask="255.255.255.240" gateway="192.35.8.161"/>
    <route hosts="routerB" destination="192.35.8.96"  netmask="255.255.255.224" gateway="192.35.8.161"/>

    <!-- ===== DEFAULT GATEWAY ROUTES  HOSTS ===== -->
    <route hosts="engHost*"       destination="0.0.0.0" netmask="0.0.0.0" gateway="192.35.8.1"/>
    <route hosts="marketingHost*" destination="0.0.0.0" netmask="0.0.0.0" gateway="192.35.8.145"/>
    <route hosts="salesHost*"     destination="0.0.0.0" netmask="0.0.0.0" gateway="192.35.8.97"/>
    <route hosts="mgmtHost*"      destination="0.0.0.0" netmask="0.0.0.0" gateway="192.35.8.65"/>
    <route hosts="adminHost*"     destination="0.0.0.0" netmask="0.0.0.0" gateway="192.35.8.129"/>

</config>
\end{lstlisting}

\subsection*{lab7.ned}
\begin{lstlisting}[style=labcode]
package lab7practice;

import inet.node.ethernet.EthernetSwitch;
import inet.node.inet.StandardHost;
import inet.node.inet.Router;
import inet.node.ethernet.Eth100M;
import inet.networklayer.configurator.ipv4.Ipv4NetworkConfigurator;

network lab7Network
{
    @display("bgb=1200,700");

    submodules:
		configurator: Ipv4NetworkConfigurator {
        	@display("p=80,50");
    	}

        // ===== ROUTERS =====
        routerA: Router {
            @display("p=450,320");
        }
        routerB: Router {
            @display("p=750,320");
        }

        // ===== ENGINEERING SUBNET (Router A) =====
        engSwitch: EthernetSwitch {
            @display("p=200,150");
        }
        engHostA: StandardHost {
            @display("p=80,80");
        }
        engHostB: StandardHost {
            @display("p=200,60");
        }
        engHostC: StandardHost {
            @display("p=320,80");
        }

        // ===== MARKETING SUBNET (Router A) =====
        marketingSwitch: EthernetSwitch {
            @display("p=200,480");
        }
        marketingHostA: StandardHost {
            @display("p=80,560");
        }
        marketingHostB: StandardHost {
            @display("p=200,600");
        }
        marketingHostC: StandardHost {
            @display("p=320,560");
        }

        // ===== SALES SUBNET (Router A) =====
        salesSwitch: EthernetSwitch {
            @display("p=572,537");
        }
        salesHostA: StandardHost {
            @display("p=486,619");
        }
        salesHostB: StandardHost {
            @display("p=617,658");
        }
        salesHostC: StandardHost {
            @display("p=691,629");
        }

        // ===== MANAGEMENT SUBNET (Router B) =====
        mgmtSwitch: EthernetSwitch {
            @display("p=1000,150");
        }
        mgmtHostA: StandardHost {
            @display("p=880,80");
        }
        mgmtHostB: StandardHost {
            @display("p=1000,60");
        }
        mgmtHostC: StandardHost {
            @display("p=1120,80");
        }

        // ===== ADMINISTRATION SUBNET (Router B) =====
        adminSwitch: EthernetSwitch {
            @display("p=1000,480");
        }
        adminHostA: StandardHost {
            @display("p=880,560");
        }
        adminHostB: StandardHost {
            @display("p=1000,600");
        }
        adminHostC: StandardHost {
            @display("p=1120,560");
        }

    connections:

        // ===== ROUTER A <--> ROUTER B (point-to-point link) =====


        // ===== ROUTER A <--> ENGINEERING SWITCH =====
        routerA.ethg++ <--> Eth100M <--> engSwitch.ethg++;
        engHostA.ethg++ <--> Eth100M <--> engSwitch.ethg++;
        engHostB.ethg++ <--> Eth100M <--> engSwitch.ethg++;
        engHostC.ethg++ <--> Eth100M <--> engSwitch.ethg++;

        // ===== ROUTER A <--> MARKETING SWITCH =====
        routerA.ethg++ <--> Eth100M <--> marketingSwitch.ethg++;
        marketingHostA.ethg++ <--> Eth100M <--> marketingSwitch.ethg++;
        marketingHostB.ethg++ <--> Eth100M <--> marketingSwitch.ethg++;
        marketingHostC.ethg++ <--> Eth100M <--> marketingSwitch.ethg++;

        // ===== ROUTER A <--> SALES SWITCH =====
        routerA.ethg++ <--> Eth100M <--> salesSwitch.ethg++;
        salesHostA.ethg++ <--> Eth100M <--> salesSwitch.ethg++;
        salesHostB.ethg++ <--> Eth100M <--> salesSwitch.ethg++;
        salesHostC.ethg++ <--> Eth100M <--> salesSwitch.ethg++;

        // ===== ROUTER B <--> MANAGEMENT SWITCH =====
        routerB.ethg++ <--> Eth100M <--> mgmtSwitch.ethg++;
        mgmtHostA.ethg++ <--> Eth100M <--> mgmtSwitch.ethg++;
        mgmtHostB.ethg++ <--> Eth100M <--> mgmtSwitch.ethg++;
        mgmtHostC.ethg++ <--> Eth100M <--> mgmtSwitch.ethg++;

        // ===== ROUTER B <--> ADMINISTRATION SWITCH =====
        routerB.ethg++ <--> Eth100M <--> adminSwitch.ethg++;
        adminHostA.ethg++ <--> Eth100M <--> adminSwitch.ethg++;
        adminHostB.ethg++ <--> Eth100M <--> adminSwitch.ethg++;
        adminHostC.ethg++ <--> Eth100M <--> adminSwitch.ethg++;

        routerA.ethg++ <--> Eth100M <--> routerB.ethg++;
}
\end{lstlisting}

\subsection*{omnetpp.ini}
\begin{lstlisting}[style=labcode]
[General]
network = lab7Network
sim-time-limit = 15s

# ===== CONFIGURATOR  address assignment only, NO auto routing =====
*.configurator.config = xmldoc("ip2.xml")
*.configurator.assignAddresses = true
*.configurator.addStaticRoutes = false
*.configurator.addDefaultRoutes = false
*.configurator.addSubnetRoutes = false


# ===== PING APPS  every HostA pings across the network =====

# Engineering HostA  pings Management and Admin (cross-router test)
*.engHostA.numApps = 2
*.engHostA.app[0].typename = "PingApp"
*.engHostA.app[0].destAddr = "salesHostA"
*.engHostA.app[0].packetSize = 500B
*.engHostA.app[0].sendInterval = 1s

*.engHostA.app[1].typename = "PingApp"
*.engHostA.app[1].destAddr = "adminHostA"
*.engHostA.app[1].packetSize = 500B
*.engHostA.app[1].sendInterval = 1s

# Marketing HostA  pings Sales HostA (same router)
*.marketingHostA.numApps = 1
*.marketingHostA.app[0].typename = "PingApp"
*.marketingHostA.app[0].destAddr = "salesHostA"
*.marketingHostA.app[0].packetSize = 500B
*.marketingHostA.app[0].sendInterval = 1s

# Sales HostA  pings Admin HostA (cross-router test)
*.salesHostA.numApps = 1
*.salesHostA.app[0].typename = "PingApp"
*.salesHostA.app[0].destAddr = "adminHostA"
*.salesHostA.app[0].packetSize = 500B
*.salesHostA.app[0].sendInterval = 1s

# Management HostA  pings Engineering HostA (cross-router test)
*.mgmtHostA.numApps = 1
*.mgmtHostA.app[0].typename = "PingApp"
*.mgmtHostA.app[0].destAddr = "engHostA"
*.mgmtHostA.app[0].packetSize = 500B
*.mgmtHostA.app[0].sendInterval = 1s

# Admin HostA  pings Marketing HostA (cross-router test)
*.adminHostA.numApps = 1
*.adminHostA.app[0].typename = "PingApp"
*.adminHostA.app[0].destAddr = "marketingHostA"
*.adminHostA.app[0].packetSize = 500B
*.adminHostA.app[0].sendInterval = 1s
\end{lstlisting}

\section{Folder: socketProg\_lab2}

\subsection*{TCP\_frame\_client.py}
\begin{lstlisting}[style=labcode]
import socket
import struct
import sys

# Use command-line argument if provided
HOST = sys.argv[1] if len(sys.argv) > 1 else '127.0.0.1'
PORT = 7007

# Create TCP socket
client = socket.socket(socket.AF_INET, socket.SOCK_STREAM)

# Connect to server
client.connect((HOST, PORT))

for msg in ["hello", "this is a longer message", "bye"]:
    # Convert string to bytes
    data = msg.encode()

    # Send: 4-byte length prefix + actual data
    client.sendall(struct.pack('!I', len(data)) + data)

    # Receive 4-byte response header
    hdr = client.recv(4)
    if not hdr:
        print("Server closed")
        break

    # Extract message length
    (length,) = struct.unpack('!I', hdr)

    # Receive full response payload
    resp = b''
    while len(resp) < length:
        chunk = client.recv(length - len(resp))
        if not chunk:
            break
        resp += chunk

    print("Echo:", resp.decode())

client.close()
\end{lstlisting}

\subsection*{TCP\_frame\_server.py}
\begin{lstlisting}[style=labcode]
import socket
import struct

# Listen on all network interfaces
HOST = '10.170.232.240'
PORT = 7007


def recv_all(sock, n):
    """
    Receive exactly n bytes from the socket.
    TCP does not guarantee that recv(n) returns n bytes,
    so we must loop until all bytes are received.
    """
    data = b''
    while len(data) < n:
        chunk = sock.recv(n - len(data))
        if not chunk:  # Connection closed
            return None
        data += chunk
    return data


# Create TCP socket (IPv4, Stream-based)
server = socket.socket(socket.AF_INET, socket.SOCK_STREAM)

# Allow quick reuse of the port after restart
server.setsockopt(socket.SOL_SOCKET, socket.SO_REUSEADDR, 1)

# Bind to address and port
server.bind((HOST, PORT))

# Listen for incoming connections (queue size 5)
server.listen(5)
print("Length-prefixed TCP server on", PORT)

# Accept one client connection
conn, addr = server.accept()
print("Connected by", addr)

try:
    while True:
        # Step 1: Read the 4-byte length header
        hdr = recv_all(conn, 4)
        if not hdr:
            break  # Client disconnected

        # '!I' = Network byte order (big-endian) unsigned int
        (length,) = struct.unpack('!I', hdr)

        # Step 2: Read exactly 'length' bytes of payload
        payload = recv_all(conn, length)
        if payload is None:
            break

        print("Got message:", payload.decode())

        # Step 3: Echo back using same framing format
        # First pack length, then append payload
        conn.sendall(struct.pack('!I', len(payload)) + payload)

finally:
    conn.close()
    server.close()
\end{lstlisting}

\section{Folder: SP\_practice}

\subsection*{client.py}
\begin{lstlisting}[style=labcode]
import socket
import threading
import sys

if len(sys.argv) != 3:
    print("Usage: python3 tcp_chat_client.py <server_ip> <port>")
    sys.exit(1)

SERVER_IP = sys.argv[1]
PORT = int(sys.argv[2])

# Create TCP socket
client = socket.socket(socket.AF_INET, socket.SOCK_STREAM)

# Connect to server
client.connect((SERVER_IP, PORT))
print("Connected to ChatIUT server. Type /quit to exit")

def receive_messages():
    while True:
        try:
            msg = client.recv(1024)
            if not msg:
                print('Server disconnected')
                break
            print(msg.decode('utf-8'))
        except:
            break

def send_messages():
    while True:
        msg = input()
        if msg.lower() == '/quit':
            client.close()
            print("Disconnected from server.")
            break
        try:
            client.sendall(msg.encode('utf-8'))
        except:
            print("Error sending message.")
            break

recv_thread = threading.Thread(target=receive_messages, daemon=True)
recv_thread.start()

send_thread = threading.Thread(target=send_messages, daemon=True)
send_thread.start()

recv_thread.join()
send_thread.join()
\end{lstlisting}

\subsection*{server.py}
\begin{lstlisting}[style=labcode]
import socket
import threading 
import logging

logging.basicConfig(
    filename='server.log',
    level=logging.INFO,
    format='%(asctime)s - %(message)s',
    datefmt='%Y-%m-%d %H:%M:%S'
)

# Create TCP socket
server = socket.socket(socket.AF_INET, socket.SOCK_STREAM)

# Bind server
server.bind(('0.0.0.0', 6006))

# Listen for connections
server.listen(5)

print("TCP server listening on port 6006...")
logging.info("TCP server listening on port 6006...")

clients = []

def broadcast(message, sender):
    for client in clients:
        if sender != client:
            try:
                client.sendall(message)
            except:
                client.remove(client)

def handle_client(client_socket, addr):
    while True:
        try:
            msg = client_socket.recv(1024)
            if not msg:
                break;
            print(f"{addr} says: {msg.decode('utf-8')}")
            logging.info(f"Forwarded message from {addr}: {msg.decode('utf-8')}")
            broadcast(msg, client_socket)
        except:
            break
    #disconnet
    logging.info(f"{addr} disconnected")
    clients.remove(client_socket)
    client_socket.close()
    print(f"{addr} disconnected")
   
# main
try:
    while True:
        client_sock, addr = server.accept()
        clients.append(client_sock)
        thread = threading.Thread(target=handle_client, args=(client_sock, addr), daemon=True)
        thread.start()
except KeyboardInterrupt:
    print("\nServer shutting down...")
    logging.info("Server stopped")
finally:
    server.close();
    server.close();
\end{lstlisting}

\end{document}

